%!TEX root = ../main.tex
% =============================================================
%  CAPÍTULO 5 — JUSTIFICACIÓN ECONÓMICA (Entrega 5 · 15 % · 15/9)
% =============================================================

\chapter{Justificación económica}\label{cap:economica}

\begin{guia}[Guía de la cátedra: qué se evalúa en este capítulo]
Es el capítulo de mayor peso individual (15\,\%). Se evalúa como una
evaluación de proyectos de inversión:
\begin{enumerate}
    \item \textbf{Análisis de mercado y competidores}: tamaño, segmentos,
          precios de referencia.
    \item \textbf{Modelo de negocio}: B2B, B2C, B2G o interno; niveles de
          precio y fuentes de ingreso.
    \item \textbf{Inversión inicial} desglosada por rol y tiempo, más
          infraestructura.
    \item \textbf{Costos operativos} recurrentes.
    \item \textbf{Escenarios} pesimista, base y optimista, con flujo de caja
          proyectado a 3 a 5 años.
    \item \textbf{Indicadores}: \gls{van}, \gls{tir} y período de repago,
          con una tasa de descuento construida y justificada con fuentes.
    \item \textbf{Supuestos y restricciones} explícitos.
\end{enumerate}
Todo número con su origen: cotización, tarifa pública, estadística
citada o estimación propia declarada como tal.
\end{guia}

\section{Alcance económico, escenarios y advertencia metodológica}

Este capítulo evalúa la factibilidad económica de MotoSafe GBA bajo tres
etapas de despliegue: demo académica (la situación actual, sin
desembolso adicional relevante), piloto público (una operación acotada
en usuarios y tiempo, orientada a validar adopción real) y operación
sostenida (un servicio disponible de forma continua, con un modelo de
ingresos que le permita cubrir sus costos). Los indicadores financieros
de este capítulo, Valor Actual Neto, Tasa Interna de Retorno y período
de recupero, solo tienen sentido para la tercera etapa: la demo
académica no requiere inversión que evaluar, y el piloto público es, por
definición, una fase de validación sin pretensión de rentabilidad
propia.

El proyecto no cuenta, a la fecha de esta entrega, con datos de mercado
reales (no se contactó a municipios ni aseguradoras para validar
disposición de pago). Siguiendo el criterio expuesto en la Clase 9 de la
cátedra para casos sin datos de mercado disponibles, este capítulo
construye un ejemplo numérico completo sobre supuestos explícitos,
declarados uno por uno y diferenciados con claridad de los hechos
verificables del proyecto (costos de infraestructura ya cotizados,
conteos de código verificados en el repositorio). Cada supuesto se marca
como tal. El objetivo no es afirmar que el proyecto generará esos
ingresos, sino demostrar el método de evaluación y su sensibilidad ante
cambios razonables en esos supuestos.

La siguiente tabla reúne, en un único lugar, todos los supuestos que
atraviesan el capítulo, para que el horizonte, la moneda y la tasa
usados en 5.2 a 5.6 se lean como un mismo conjunto coherente y no como
cifras sueltas por sección.

\begin{longtable}[]{@{}
  >{\raggedright\arraybackslash}p{(\columnwidth - 4\tabcolsep) * \real{0.2766}}
  >{\raggedright\arraybackslash}p{(\columnwidth - 4\tabcolsep) * \real{0.2553}}
  >{\raggedright\arraybackslash}p{(\columnwidth - 4\tabcolsep) * \real{0.4681}}@{}}
\caption{Supuestos de modelado y hechos verificables que atraviesan el Capítulo 5.}\label{tab:5-1}\\
\toprule\noalign{}
\begin{minipage}[b]{\linewidth}\raggedright
\textbf{Supuesto}
\end{minipage} & \begin{minipage}[b]{\linewidth}\raggedright
\textbf{Valor adoptado}
\end{minipage} & \begin{minipage}[b]{\linewidth}\raggedright
\textbf{Naturaleza y justificación}
\end{minipage} \\
\midrule\noalign{}
\endfirsthead
\toprule\noalign{}
\begin{minipage}[b]{\linewidth}\raggedright
\textbf{Supuesto}
\end{minipage} & \begin{minipage}[b]{\linewidth}\raggedright
\textbf{Valor adoptado}
\end{minipage} & \begin{minipage}[b]{\linewidth}\raggedright
\textbf{Naturaleza y justificación}
\end{minipage} \\
\midrule\noalign{}
\endhead
\bottomrule\noalign{}
\endlastfoot
Horizonte de evaluación & 5 años (Año 0 a Año 5) & Supuesto: horizonte
típico para un proyecto de servicio digital sin activos físicos,
suficiente para observar si los ingresos proyectados alcanzan a cubrir
los costos operativos. \\
Moneda de todo el flujo & Dólares estadounidenses (USD) & Supuesto de
modelado: evita proyectar una tasa de inflación local y es consistente
con la cotización en USD de los servicios de infraestructura (Render,
TomTom). \\
Depreciación de activos & No se cargan cargos por depreciación & Hecho
conceptual: el proyecto es un servicio digital sin inventario físico
relevante que amortizar. \\
Tasa de descuento & 17\% anual (16,65\% calculado, redondeado) &
Supuesto compuesto: costo de oportunidad (bono EE. UU. a 10 años, hecho
verificado, \citealp{cnbc2026}) + riesgo país (EMBI+ Argentina, hecho
verificado, \citealp{rava2026}) + prima de riesgo propia del proyecto
de 7 puntos (supuesto). \\
Crecimiento anual de costos operativos & 10\% anual sobre el año
anterior & Supuesto: refleja el crecimiento esperado de tráfico, volumen
de prensa almacenada y consumo de TomTom a medida que crece la base de
usuarios. \\
Escenario de ingresos (licencia B2B) & USD 0 / 400 / 1.200 / 2.200 /
3.500 (años 1 a 5) & Supuesto de modelado, sin contacto comercial
iniciado con municipios ni aseguradoras a la fecha de esta entrega. \\
Suscripción Premium (usuarios finales) & USD 10/año; 0 / 100 / 250 / 400
/ 550 usuarios pagos (años 1 a 5) & Supuesto de modelado, segunda fuente
de ingresos sin base de usuarios validada a la fecha de esta entrega. \\
Publicidad (plan gratuito) & CPM USD 0,50; 1 ruta iniciada/semana por
usuario; conversión a Premium del 2\% sobre la base total & Supuesto de
modelado, tercera fuente de ingresos, sin anunciante contactado a la
fecha de esta entrega. \\
Inversión inicial (Año 0) & USD 5.740 & Supuesto: incluye el desarrollo
ya realizado del MVP y las tareas adicionales para escalarlo a operación
sostenida (ver 5.2). \\
Costo de infraestructura verificado & USD 0/mes (demo/piloto actual:
planes Free de Render y Supabase); USD 67,65/mes (Render Starter +
Supabase Pro + Open-Meteo en su tier comercial, con el frontend movido a
hosting estático gratuito) para operación sostenida & Hecho verificable
en ambos extremos: el gratuito corre hoy en producción; el pago está
cotizado directamente en cada proveedor (Render y Supabase). \\
\end{longtable}

\section{Inversión inicial (etapa de operación sostenida)}

La inversión inicial reúne, en un único total, el desarrollo ya
realizado del MVP (backend, frontend, modelo y pipeline de datos) y las
tareas adicionales necesarias para pasar de ese MVP de demostración a un
servicio público sostenido. El desarrollo ya realizado incluye tanto las
horas propias como el trabajo de los asistentes de inteligencia
artificial usados bajo esa supervisión, valorizado a la misma tarifa de
referencia usada en todo este capítulo (USD 12/hora, desarrollador
junior en Argentina, supuesto propio). Las tareas adicionales
corresponden a: endurecimiento de seguridad antes de exponer el sistema
(migración de tokens de borrado existentes y revisión de dependencias;
la cuota de envío de reportes ya está implementada), puesta en marcha de
infraestructura productiva (base de datos persistente, dominio propio,
TLS, backups automatizados) y una campaña mínima de difusión inicial en
comunidades de motociclistas del AMBA.

\begin{longtable}[]{@{}
  >{\raggedright\arraybackslash}p{(\columnwidth - 6\tabcolsep) * \real{0.3696}}
  >{\raggedright\arraybackslash}p{(\columnwidth - 6\tabcolsep) * \real{0.1739}}
  >{\raggedright\arraybackslash}p{(\columnwidth - 6\tabcolsep) * \real{0.1957}}
  >{\raggedright\arraybackslash}p{(\columnwidth - 6\tabcolsep) * \real{0.2609}}@{}}
\caption{Desglose de la inversión inicial por concepto, horas y naturaleza (supuesto o hecho).}\label{tab:5-2}\\
\toprule\noalign{}
\begin{minipage}[b]{\linewidth}\raggedright
\textbf{Concepto}
\end{minipage} & \begin{minipage}[b]{\linewidth}\raggedright
\textbf{Horas}
\end{minipage} & \begin{minipage}[b]{\linewidth}\raggedright
\textbf{Costo (USD)}
\end{minipage} & \begin{minipage}[b]{\linewidth}\raggedright
\textbf{Naturaleza}
\end{minipage} \\
\midrule\noalign{}
\endfirsthead
\toprule\noalign{}
\begin{minipage}[b]{\linewidth}\raggedright
\textbf{Concepto}
\end{minipage} & \begin{minipage}[b]{\linewidth}\raggedright
\textbf{Horas}
\end{minipage} & \begin{minipage}[b]{\linewidth}\raggedright
\textbf{Costo (USD)}
\end{minipage} & \begin{minipage}[b]{\linewidth}\raggedright
\textbf{Naturaleza}
\end{minipage} \\
\midrule\noalign{}
\endhead
\bottomrule\noalign{}
\endlastfoot
Investigación y definición del proyecto & 80 & 960 & Supuesto: horas
estimadas según el alcance real del trabajo ya hecho, sin registro
cronometrado \\
Desarrollo e integraciones (backend, modelo, frontend, scraper de
prensa) & 220 & 2.640 & Supuesto: horas estimadas según el alcance real
del trabajo ya hecho, sin registro cronometrado \\
Auditoría técnica, documentación y redacción académica & 40 & 480 &
Supuesto: horas estimadas según el alcance real del trabajo ya hecho,
sin registro cronometrado \\
Suscripciones de las herramientas de inteligencia artificial utilizadas
(Claude y ChatGPT/Codex) & --- & 160 & Hecho verificable en el precio
(USD 20/mes cada suscripción) y en la duración real de uso hasta la
fecha de esta entrega: primer mes (mayo de 2026) con Claude y
ChatGPT/Codex a la vez (USD 40), de junio a noviembre de 2026 (seis
meses) solo Claude (USD 20/mes); total USD 160 sobre siete meses de
uso \\
Endurecimiento de seguridad y auditoría previa a escalar & --- & 800 &
Supuesto: horas de desarrollo valorizadas \\
Infraestructura productiva inicial (BD persistente, dominio, TLS,
backups) & --- & 300 & Supuesto: cotización estimada sobre tarifas de
Render, Supabase y servicios de dominio/TLS \\
Difusión inicial en comunidades de motociclistas & --- & 400 & Supuesto:
presupuesto de lanzamiento \\
\textbf{Inversión inicial total (Año 0)} & \textbf{340} & \textbf{5.740}
& \textbf{Supuesto} \\
\end{longtable}

Las horas del desarrollo ya realizado se estiman por analogía con
proyectos de alcance e integración comparable (backend propio, modelo de
aprendizaje automático, frontend geoespacial y proceso de ingesta
programado). El costo de las suscripciones de inteligencia artificial
usa el precio real de cada plan individual (USD 20/mes Claude, USD
20/mes ChatGPT/Codex) y la duración real de uso hasta la fecha de esta
entrega, no un supuesto: USD 40 en mayo de 2026, el único mes en que se
usaron ambas herramientas, y USD 20/mes desde junio hasta noviembre de
2026 (seis meses), en que solo se mantuvo Claude; total USD 160 sobre
siete meses. Todas las horas y esta tarifa son un supuesto propio,
declarado como tal, y no una cotización ni un registro de horas
cronometradas.

\section{Costos operativos}

La arquitectura real, verificada en producción al cierre de esta
entrega, corre en su totalidad en planes gratuitos: Render (API y sitio
web, Web Services sin costo), Supabase (PostgreSQL con PostGIS, plan
Free sin vencimiento) y GitHub Actions (scraper de prensa por cron). Ese
hecho sostiene las etapas de demo académica y piloto público a costo
cero de infraestructura. La etapa de operación sostenida, en cambio,
necesita planes pagos para evitar las limitaciones de los niveles
gratuitos (el servicio web gratuito de Render se duerme sin tráfico; el
proyecto gratuito de Supabase se pausa tras una semana sin actividad).
Esta sección proyecta esa etapa con un ajuste que reduce el costo sin
cambiar de proveedor de base de datos: el frontend es una aplicación de
una sola página sin lógica de servidor, así que para esa etapa se
proyecta publicarlo en un hosting estático gratuito (p. ej. Cloudflare
Pages), dejando en Render pago solo la API (Starter, USD 7/mes);
Supabase se mantiene como base de datos, en su plan Pro (USD 25/mes)
para esa misma etapa. Ese cambio no es solo una migración de hosting:
hoy el frontend y la API conviven detrás del mismo proxy nginx, y mover
el frontend a un dominio estático separado exige reconfigurar CORS y
pasar la cookie de sesión a un esquema entre dominios (SameSite=None,
Secure), trabajo de implementación todavía no hecho, no un simple cambio
de proveedor. El clima en vivo cambia de categoría en esta misma etapa:
hoy consulta la API gratuita y no comercial de Open-Meteo, pero su
política reserva ese nivel gratuito a uso no comercial, así que una
operación que cobra por licencias de datos o suscripciones necesita su
tier comercial (aproximadamente USD 35,65/mes). La siguiente tabla
identifica los rubros de costo variable y el factor que más los impacta;
los montos anuales de la etapa de operación sostenida (columna derecha)
son un supuesto de dimensionamiento sobre esas tarifas, no una
cotización cerrada por volumen.

\begin{longtable}[]{@{}
  >{\raggedright\arraybackslash}p{(\columnwidth - 6\tabcolsep) * \real{0.3171}}
  >{\raggedright\arraybackslash}p{(\columnwidth - 6\tabcolsep) * \real{0.3171}}
  >{\raggedright\arraybackslash}p{(\columnwidth - 6\tabcolsep) * \real{0.2195}}
  >{\raggedright\arraybackslash}p{(\columnwidth - 6\tabcolsep) * \real{0.1463}}@{}}
\caption{Costos operativos proyectados por rubro y variable de mayor impacto.}\label{tab:5-3}\\
\toprule\noalign{}
\begin{minipage}[b]{\linewidth}\raggedright
\textbf{Rubro}
\end{minipage} & \begin{minipage}[b]{\linewidth}\raggedright
\textbf{Variable que más impacta el costo}
\end{minipage} & \begin{minipage}[b]{\linewidth}\raggedright
\textbf{Base (hecho)}
\end{minipage} & \begin{minipage}[b]{\linewidth}\raggedright
\textbf{Año 1 estimado (supuesto)}
\end{minipage} \\
\midrule\noalign{}
\endfirsthead
\toprule\noalign{}
\begin{minipage}[b]{\linewidth}\raggedright
\textbf{Rubro}
\end{minipage} & \begin{minipage}[b]{\linewidth}\raggedright
\textbf{Variable que más impacta el costo}
\end{minipage} & \begin{minipage}[b]{\linewidth}\raggedright
\textbf{Base (hecho)}
\end{minipage} & \begin{minipage}[b]{\linewidth}\raggedright
\textbf{Año 1 estimado (supuesto)}
\end{minipage} \\
\midrule\noalign{}
\endhead
\bottomrule\noalign{}
\endlastfoot
Render (solo API) & Requests, RAM/CPU activa & USD 0/mes verificado
(Free, demo actual); USD 7/mes (Starter, sin sleep) & USD 84 \\
Frontend (Cloudflare Pages) & Visitas, ancho de banda & USD 0/mes
verificado en demo (hoy en Render Free); se proyecta mover a hosting
estático gratuito para operación sostenida, sin costo dentro de los
límites gratuitos del proveedor (sujeto a sus topes de uso) & USD 0 \\
Supabase (PostgreSQL/PostGIS) & Tamaño de datos, cómputo, usuarios
activos & USD 0/mes verificado (Free, sin vencimiento, demo actual); USD
25/mes (Pro) & USD 300 \\
Mapas y ruteo (TomTom) & Búsquedas, rutas, consulta de tráfico & Cuota
gratuita limitada & USD 180 \\
Scraper de prensa (GitHub Actions) & Minutos de cómputo por corrida (hoy
cada 3 h) & USD 0/mes verificado, incluso a esta escala (cuota de
minutos muy por encima del uso real) & USD 0 \\
Envío de mail (API de Gmail) & Volumen de envíos & USD 0/mes verificado,
cuota muy por encima del uso necesario & USD 0 \\
Clima en vivo (Open-Meteo) & Llamadas por request, uso comercial & USD
0/mes verificado en demo (tier no comercial); requiere su tier comercial
(\textasciitilde USD 35,65/mes) para operación sostenida con licencias
pagas & USD 428 \\
Dominio, TLS y observabilidad & Proveedor y retención de logs & Costo
anual estándar (pendiente de contratar) & USD 56 \\
Costo operativo total Año 1 (demo verificado: USD 0; operación
sostenida, estimado) & & & USD 1.048 \\
\end{longtable}

El costo operativo total se proyecta con un crecimiento del 10\% anual
sobre el año anterior (supuesto), reflejando el aumento esperado de
tráfico, volumen de prensa almacenada y consumo de TomTom y Open-Meteo a
medida que crece la base de usuarios: USD 1.048 (año 1), USD 1.153 (año
2), USD 1.268 (año 3), USD 1.395 (año 4) y USD 1.534 (año 5).

\section{Escenario de ingresos}

SUPUESTO (no validado con el mercado): el escenario de operación
sostenida financia sus costos mediante tres fuentes combinadas, ninguna
con contacto comercial iniciado a la fecha de esta entrega. La primera
es una licencia de datos agregados y anonimizados de riesgo vial (capas
de hotspots, no datos personales) ofrecida a municipios o aseguradoras
como servicio de suscripción a partir del tercer año, una vez que el
volumen de datos históricos lo justifique. La segunda es una suscripción
Premium opcional para usuarios finales, a USD 10 por año, con un
diferencial central: viaje compartido en vivo, donde el usuario comparte
su ubicación y ruta en tiempo real con un contacto de confianza mientras
maneja. Sobre eso se suman, como extras menores, pantalla sin
publicidad, guardado ilimitado de rutas favoritas, alertas
personalizadas por zona y acceso sin conexión a las capas de riesgo. El
plan gratuito conserva todo lo que el sistema ya tiene implementado hoy
(mapa con las cuatro capas de riesgo, comparador de rutas, reportes
comunitarios, panel de prensa, hasta 3 rutas favoritas guardadas), pero
no incluye el viaje compartido en vivo ni los extras mencionados. La
tercera fuente es publicidad mostrada solo a usuarios del plan gratuito,
en la pantalla de carga antes de iniciar una ruta, momento en el que se
proyecta que el usuario todavía esté detenido, sin garantía de que lo
esté siempre; esa suposición de seguridad queda sujeta a evaluación
antes de implementarse, no es un hecho verificado. El plan Premium no
muestra publicidad. Ninguna de las tres fuentes es una funcionalidad ya
construida en el sistema, ni hay una negociación comercial en curso con
un anunciante, con municipios o con aseguradoras: las tres son supuestos
de modelado, no un producto ya ofrecido.

Para la licencia de datos se proyecta: USD 0 (año 1, etapa de
construcción de base de usuarios sin monetizar), USD 400 (año 2, un
auspicio puntual de una organización interesada, distinto de una
licencia de datos formal pero agrupado en esta misma fuente por
simplicidad), USD 1.200 (año 3, primera licencia piloto de datos ya
formal), USD 2.200 (año 4) y USD 3.500 (año 5). Para la suscripción
Premium, a partir del año 2 se supone una conversión del 2\% (supuesto,
conversión freemium típica de una app de consumo) sobre una base de
usuarios activos totales estimada en 5.000 (año 2), 12.500 (año 3),
20.000 (año 4) y 27.500 (año 5); el año 1 se proyecta directamente en 0
usuarios pagos, no por la fórmula del 2\% (que daría 40 sobre sus 2.000
usuarios activos estimados), sino porque esa primera etapa se supone
dedicada a construir base de usuarios antes de ofrecer el plan pago. Eso
da 0 usuarios pagos en el año 1, 100 en el año 2, 250 en el año 3, 400
en el año 4 y 550 en el año 5. Para la publicidad, mostrada a los
usuarios gratuitos de esa misma base (la base total menos los usuarios
Premium), se adoptó el extremo conservador de dos variables sin validar:
una ruta iniciada por semana por usuario activo, y un CPM (costo por
cada 1.000 vistas) de USD 0,50, en la parte baja del rango típico de
publicidad móvil en Latinoamérica. Eso proyecta USD 52 (año 1), USD 127
(año 2), USD 319 (año 3), USD 510 (año 4) y USD 701 (año 5).

\begin{longtable}[]{@{}
  >{\raggedright\arraybackslash}p{(\columnwidth - 8\tabcolsep) * \real{0.1513}}
  >{\raggedright\arraybackslash}p{(\columnwidth - 8\tabcolsep) * \real{0.2089}}
  >{\raggedright\arraybackslash}p{(\columnwidth - 8\tabcolsep) * \real{0.2129}}
  >{\raggedright\arraybackslash}p{(\columnwidth - 8\tabcolsep) * \real{0.2194}}
  >{\raggedright\arraybackslash}p{(\columnwidth - 8\tabcolsep) * \real{0.2076}}@{}}
\caption{Proyección de ingresos por licencia B2B, suscripción Premium y publicidad, años 1 a 5.}\label{tab:5-4}\\
\toprule\noalign{}
\begin{minipage}[b]{\linewidth}\raggedright
\textbf{Año}
\end{minipage} & \begin{minipage}[b]{\linewidth}\raggedright
\textbf{Licencia B2B (USD)}
\end{minipage} & \begin{minipage}[b]{\linewidth}\raggedright
\textbf{Premium (USD)}
\end{minipage} & \begin{minipage}[b]{\linewidth}\raggedright
\textbf{Publicidad (USD)}
\end{minipage} & \begin{minipage}[b]{\linewidth}\raggedright
\textbf{Total ingresos (USD)}
\end{minipage} \\
\midrule\noalign{}
\endfirsthead
\toprule\noalign{}
\begin{minipage}[b]{\linewidth}\raggedright
\textbf{Año}
\end{minipage} & \begin{minipage}[b]{\linewidth}\raggedright
\textbf{Licencia B2B (USD)}
\end{minipage} & \begin{minipage}[b]{\linewidth}\raggedright
\textbf{Premium (USD)}
\end{minipage} & \begin{minipage}[b]{\linewidth}\raggedright
\textbf{Publicidad (USD)}
\end{minipage} & \begin{minipage}[b]{\linewidth}\raggedright
\textbf{Total ingresos (USD)}
\end{minipage} \\
\midrule\noalign{}
\endhead
\bottomrule\noalign{}
\endlastfoot
1 & 0 & 0 & 52 & 52 \\
2 & 400 & 1.000 & 127 & 1.527 \\
3 & 1.200 & 2.500 & 319 & 4.019 \\
4 & 2.200 & 4.000 & 510 & 6.710 \\
5 & 3.500 & 5.500 & 701 & 9.701 \\
\end{longtable}

\section{Flujo de fondos, tasa de descuento, VAN, TIR y período de recupero}

El flujo de fondos neto de cada año surge de restar el costo operativo
(5.3) al ingreso proyectado (5.4), sin cargos por depreciación de
activos fijos por tratarse de un servicio digital sin inventario físico
relevante.

\begin{longtable}[]{@{}
  >{\raggedright\arraybackslash}p{(\columnwidth - 8\tabcolsep) * \real{0.2051}}
  >{\raggedright\arraybackslash}p{(\columnwidth - 8\tabcolsep) * \real{0.1923}}
  >{\raggedright\arraybackslash}p{(\columnwidth - 8\tabcolsep) * \real{0.1923}}
  >{\raggedright\arraybackslash}p{(\columnwidth - 8\tabcolsep) * \real{0.2051}}
  >{\raggedright\arraybackslash}p{(\columnwidth - 8\tabcolsep) * \real{0.2051}}@{}}
\caption{Flujo de fondos neto y acumulado, años 0 a 5.}\label{tab:5-5}\\
\toprule\noalign{}
\begin{minipage}[b]{\linewidth}\raggedright
\textbf{Año}
\end{minipage} & \begin{minipage}[b]{\linewidth}\raggedright
\textbf{Ingreso (USD)}
\end{minipage} & \begin{minipage}[b]{\linewidth}\raggedright
\textbf{Costo (USD)}
\end{minipage} & \begin{minipage}[b]{\linewidth}\raggedright
\textbf{Flujo neto (USD)}
\end{minipage} & \begin{minipage}[b]{\linewidth}\raggedright
\textbf{Flujo acumulado (USD)}
\end{minipage} \\
\midrule\noalign{}
\endfirsthead
\toprule\noalign{}
\begin{minipage}[b]{\linewidth}\raggedright
\textbf{Año}
\end{minipage} & \begin{minipage}[b]{\linewidth}\raggedright
\textbf{Ingreso (USD)}
\end{minipage} & \begin{minipage}[b]{\linewidth}\raggedright
\textbf{Costo (USD)}
\end{minipage} & \begin{minipage}[b]{\linewidth}\raggedright
\textbf{Flujo neto (USD)}
\end{minipage} & \begin{minipage}[b]{\linewidth}\raggedright
\textbf{Flujo acumulado (USD)}
\end{minipage} \\
\midrule\noalign{}
\endhead
\bottomrule\noalign{}
\endlastfoot
0 & --- & 5.740 (inversión) & -5.740 & -5.740 \\
1 & 52 & 1.048 & -996 & -6.736 \\
2 & 1.527 & 1.153 & 374 & -6.362 \\
3 & 4.019 & 1.268 & 2.751 & -3.611 \\
4 & 6.710 & 1.395 & 5.315 & 1.704 \\
5 & 9.701 & 1.534 & 8.167 & 9.871 \\
\end{longtable}

Tasa de descuento (supuesto, construido en dólares estadounidenses según
el método de tres componentes de la Clase 9): costo de oportunidad,
aproximado por el rendimiento del bono del Tesoro de Estados Unidos a 10
años, que cotizó en torno a 4,8\% \citep{cnbc2026}; protección contra
inflación, incorporada de forma implícita al modelar el flujo
íntegramente en dólares en lugar de en pesos argentinos, evitando así
proyectar una tasa de inflación local; y compensación de riesgo,
compuesta por el riesgo país argentino (EMBI+ Argentina, 485 puntos
básicos según Rava Bursátil) más una prima adicional de riesgo
específico del proyecto de 7 puntos porcentuales (supuesto, que refleja
la ausencia de historial de ingresos y la dependencia de fuentes de
ingreso todavía no contactadas comercialmente). La suma de los tres
componentes (4,8\% + 4,85\% + 7\%) da una tasa de descuento de 16,65\%,
redondeada a 17\% para el cálculo.

VAN al 17\%: -5.740 + (-996/1,17) + (374/1,17²) + (2.751/1,17³) +
(5.315/1,17⁴) + (8.167/1,17⁵) = -5.740 + (-851,28) + 273,21 + 1.717,64 +
2.836,35 + 3.725,06 ≈ +USD 1.961. TIR: por bisección, la tasa que lleva
el VAN a cero es aproximadamente 24,8\%, con un margen de 7,8 puntos
sobre la tasa de corte del 17\%: sumar publicidad al plan gratuito
(conservadora: una ruta iniciada por semana, CPM de USD 0,50) sobre la
suscripción Premium de USD 10 por año casi duplica el margen que dejaba
la sola suba de precio de Premium. Período de recupero simple: el flujo
acumulado pasa de negativo a positivo entre el año 3 (-3.611) y el año 4
(+5.315 de flujo neto ese año), a los 3.611/5.315 = 0,68 del año, es
decir, 3 años y aproximadamente 8 meses. Período de recupero descontado:
sobre los flujos descontados acumulados (-5.740, -6.591,28, -6.318,07,
-4.600,43, -1.764,08, +1.960,98), el cruce ocurre entre el año 4 y el
año 5, a 1.764,08/3.725,06 = 0,47 del año, es decir, 4 años y
aproximadamente 6 meses.

\begin{longtable}[]{@{}
  >{\raggedright\arraybackslash}p{(\columnwidth - 6\tabcolsep) * \real{0.2020}}
  >{\raggedright\arraybackslash}p{(\columnwidth - 6\tabcolsep) * \real{0.2525}}
  >{\raggedright\arraybackslash}p{(\columnwidth - 6\tabcolsep) * \real{0.2525}}
  >{\raggedright\arraybackslash}p{(\columnwidth - 6\tabcolsep) * \real{0.2929}}@{}}
\caption{Resultado de los indicadores financieros frente a su criterio de aceptación.}\label{tab:5-6}\\
\toprule\noalign{}
\begin{minipage}[b]{\linewidth}\raggedright
\textbf{Indicador}
\end{minipage} & \begin{minipage}[b]{\linewidth}\raggedright
\textbf{Resultado}
\end{minipage} & \begin{minipage}[b]{\linewidth}\raggedright
\textbf{Criterio}
\end{minipage} & \begin{minipage}[b]{\linewidth}\raggedright
\textbf{Conclusión}
\end{minipage} \\
\midrule\noalign{}
\endfirsthead
\toprule\noalign{}
\begin{minipage}[b]{\linewidth}\raggedright
\textbf{Indicador}
\end{minipage} & \begin{minipage}[b]{\linewidth}\raggedright
\textbf{Resultado}
\end{minipage} & \begin{minipage}[b]{\linewidth}\raggedright
\textbf{Criterio}
\end{minipage} & \begin{minipage}[b]{\linewidth}\raggedright
\textbf{Conclusión}
\end{minipage} \\
\midrule\noalign{}
\endhead
\bottomrule\noalign{}
\endlastfoot
VAN (17\%) & ≈ +USD 1.961 & Aceptar si VAN ≥ 0 & Se acepta, margen
sensible a la adopción de ingresos (ver 5.6) \\
TIR & ≈ 24,8\% & Aceptar si tasa de corte ≤ TIR & Se acepta, margen de
7,8 puntos \\
Recupero simple & ≈ 3 años y 8 meses & Cuanto menor, mejor liquidez &
Recupero dentro del horizonte evaluado \\
Recupero descontado & ≈ 4 años y 6 meses & Cuanto menor, mejor liquidez
real & Recupero dentro del horizonte evaluado \\
\end{longtable}

Bajo los supuestos declarados, incluyendo el costo completo del
desarrollo ya realizado, la cotización real de infraestructura pagada
(Render Starter para la API, frontend en hosting estático gratuito,
Supabase Pro y el tier comercial de Open-Meteo), la suscripción Premium
de USD 10 por año y la publicidad en el plan gratuito propuestas en 5.4,
el escenario de operación sostenida resulta económicamente viable, con
un margen mayor al estimado en una cotización anterior de este mismo
capítulo: un VAN de USD 1.961 y una TIR que supera la tasa de corte por
7,8 puntos porcentuales. La tercera fuente de ingresos (publicidad) suma
un margen de seguridad frente al costo del clima comercial, que las
primeras dos fuentes por sí solas cubrían de forma más ajustada; esa
dependencia de tres fuentes de ingresos combinadas, ninguna validada
todavía con el mercado, se pone a prueba en el punto 5.6.

\section{Análisis de sensibilidad}

Se recalculó el VAN ante dos escenarios de desvío de una sola variable a
la vez respecto del caso base, manteniendo la tasa de descuento en 17\%.
Un sobrecosto operativo uniforme del 20\% sobre los valores de 5.3 en
todos los años lleva el VAN a aproximadamente +USD 1.166: el proyecto se
mantiene viable, con margen todavía razonable. Una adopción de ingresos
(licencia B2B, Premium y publicidad combinadas) un 20\% más lenta que la
proyectada en 5.4 lleva el VAN a aproximadamente -USD 374: el proyecto
deja de ser viable, aunque por un margen mucho más chico que antes de
sumar la publicidad. El resultado resiste mejor un sobrecosto operativo
que una demora en la adopción de ingresos, que sigue siendo la variable
más sensible del capítulo.

\begin{longtable}[]{@{}
  >{\raggedright\arraybackslash}p{(\columnwidth - 4\tabcolsep) * \real{0.3333}}
  >{\raggedright\arraybackslash}p{(\columnwidth - 4\tabcolsep) * \real{0.3333}}
  >{\raggedright\arraybackslash}p{(\columnwidth - 4\tabcolsep) * \real{0.3333}}@{}}
\caption{Sensibilidad del VAN ante el desvío de una variable a la vez.}\label{tab:5-7}\\
\toprule\noalign{}
\begin{minipage}[b]{\linewidth}\raggedright
\textbf{Desvío de una variable}
\end{minipage} & \begin{minipage}[b]{\linewidth}\raggedright
\textbf{VAN resultante (USD)}
\end{minipage} & \begin{minipage}[b]{\linewidth}\raggedright
\textbf{Lectura}
\end{minipage} \\
\midrule\noalign{}
\endfirsthead
\toprule\noalign{}
\begin{minipage}[b]{\linewidth}\raggedright
\textbf{Desvío de una variable}
\end{minipage} & \begin{minipage}[b]{\linewidth}\raggedright
\textbf{VAN resultante (USD)}
\end{minipage} & \begin{minipage}[b]{\linewidth}\raggedright
\textbf{Lectura}
\end{minipage} \\
\midrule\noalign{}
\endhead
\bottomrule\noalign{}
\endlastfoot
Caso base (5.5) & ≈ +1.961 & Viable \\
Costos operativos +20\% & ≈ +1.166 & Viable, margen amplio \\
Ingresos 20\% más lentos & ≈ -374 & No viable, pero por un margen
chico \\
\end{longtable}

Además del desvío de una variable a la vez, se construyeron tres
escenarios combinados (pesimista, base y optimista) para leer el rango
completo de resultados posibles bajo los mismos supuestos de horizonte,
moneda y tasa de descuento de la tabla de 5.1. El escenario optimista
supone una adopción de ingresos un 20\% más rápida que la proyectada en
5.4, con los costos operativos de 5.3 sin cambios. El escenario
pesimista combina, en simultáneo, el sobrecosto operativo del 20\% y la
adopción de ingresos un 20\% más lenta ya calculados por separado en la
tabla anterior; no es un tercer supuesto nuevo, sino la superposición de
los dos desvíos individuales sobre el mismo horizonte de 5 años.

\begin{longtable}[]{@{}
  >{\raggedright\arraybackslash}p{(\columnwidth - 6\tabcolsep) * \real{0.2525}}
  >{\raggedright\arraybackslash}p{(\columnwidth - 6\tabcolsep) * \real{0.2525}}
  >{\raggedright\arraybackslash}p{(\columnwidth - 6\tabcolsep) * \real{0.2525}}
  >{\raggedright\arraybackslash}p{(\columnwidth - 6\tabcolsep) * \real{0.2424}}@{}}
\caption{Sensibilidad del VAN ante escenarios combinados (optimista, base y pesimista).}\label{tab:5-8}\\
\toprule\noalign{}
\begin{minipage}[b]{\linewidth}\raggedright
\textbf{Escenario}
\end{minipage} & \begin{minipage}[b]{\linewidth}\raggedright
\textbf{Ingresos vs. base}
\end{minipage} & \begin{minipage}[b]{\linewidth}\raggedright
\textbf{Costos vs. base}
\end{minipage} & \begin{minipage}[b]{\linewidth}\raggedright
\textbf{VAN resultante (USD)}
\end{minipage} \\
\midrule\noalign{}
\endfirsthead
\toprule\noalign{}
\begin{minipage}[b]{\linewidth}\raggedright
\textbf{Escenario}
\end{minipage} & \begin{minipage}[b]{\linewidth}\raggedright
\textbf{Ingresos vs. base}
\end{minipage} & \begin{minipage}[b]{\linewidth}\raggedright
\textbf{Costos vs. base}
\end{minipage} & \begin{minipage}[b]{\linewidth}\raggedright
\textbf{VAN resultante (USD)}
\end{minipage} \\
\midrule\noalign{}
\endhead
\bottomrule\noalign{}
\endlastfoot
Optimista & +20\% (adopción más rápida) & Sin cambios & ≈ +4.296 \\
Base (5.5) & Proyección de 5.4 & Proyección de 5.3 & ≈ +1.961 \\
Pesimista & -20\% (adopción más lenta) & +20\% (sobrecosto combinado) &
≈ -1.169 \\
\end{longtable}

El rango entre el escenario optimista (+USD 4.296) y el pesimista (-USD
1.169) muestra que, con las tres fuentes de ingreso combinadas, el caso
base queda más lejos de revertirse que en la versión anterior del
capítulo: el resultado sigue dependiendo de que la adopción de ingresos
no se retrase de forma simultánea y considerable, pero la publicidad le
da un colchón adicional que las otras dos fuentes por sí solas no
tenían.

Bajo los supuestos declarados en este capítulo, incluyendo la
suscripción Premium de USD 10 por año y la publicidad en el plan
gratuito propuestas en 5.4, la operación sostenida de MotoSafe GBA
resulta económicamente viable, con un margen razonable pero condicionado
a que las tres fuentes de ingresos proyectadas (licencia de datos,
suscripción Premium y publicidad), combinadas, se validen con evidencia
real de mercado antes de escalar: el análisis de sensibilidad solo probó
una adopción de ingresos combinados hasta un 20\% más lenta que la
proyectada, no la pérdida completa de alguna de las tres fuentes, un
escenario que este capítulo no cuantifica. La decisión de migrar de la
etapa de piloto público a la de operación sostenida no debería tomarse
únicamente con este ejemplo numérico: debe estar precedida por una
validación real de la disposición a pagar (ya sea de motociclistas por
el plan Premium o de municipios y aseguradoras por la licencia de datos)
y por el interés real de un anunciante en pautar dentro de la app.
Mientras esa validación no exista, la recomendación es sostener el
proyecto en la etapa de piloto público, con límites de gasto activos en
todos los servicios externos, y recién evaluar la migración a operación
sostenida cuando exista evidencia real de al menos un trimestre de uso,
de conversión a la suscripción Premium y de interés publicitario
efectivo.

\section{Riesgos económicos operativos}

Además de la sensibilidad financiera del punto 5.6, quedan identificados
tres riesgos operativos con impacto económico directo, ya señalados en
el Capítulo 4.

\begin{longtable}[]{@{}
  >{\raggedright\arraybackslash}p{(\columnwidth - 4\tabcolsep) * \real{0.3232}}
  >{\raggedright\arraybackslash}p{(\columnwidth - 4\tabcolsep) * \real{0.2222}}
  >{\raggedright\arraybackslash}p{(\columnwidth - 4\tabcolsep) * \real{0.4545}}@{}}
\caption{Riesgos operativos con impacto económico directo y su mitigación.}\label{tab:5-9}\\
\toprule\noalign{}
\begin{minipage}[b]{\linewidth}\raggedright
\textbf{Riesgo}
\end{minipage} & \begin{minipage}[b]{\linewidth}\raggedright
\textbf{Probabilidad / impacto}
\end{minipage} & \begin{minipage}[b]{\linewidth}\raggedright
\textbf{Mitigación}
\end{minipage} \\
\midrule\noalign{}
\endfirsthead
\toprule\noalign{}
\begin{minipage}[b]{\linewidth}\raggedright
\textbf{Riesgo}
\end{minipage} & \begin{minipage}[b]{\linewidth}\raggedright
\textbf{Probabilidad / impacto}
\end{minipage} & \begin{minipage}[b]{\linewidth}\raggedright
\textbf{Mitigación}
\end{minipage} \\
\midrule\noalign{}
\endhead
\bottomrule\noalign{}
\endlastfoot
Pausado del proyecto gratuito de Supabase tras una semana sin actividad
& Alta / medio & Reactivar antes de cada demostración, o migrar a un
plan pago (Supabase Pro) antes de sostener una operación continua \\
Consumo de la API de TomTom sin límites de gasto configurados & Media /
alto & Configurar límites de gasto y monitorear el consumo real antes de
escalar el tráfico \\
Crecimiento del volumen de prensa almacenada al ampliar fuentes & Media
/ bajo a mediano plazo & Revisar la política de retención y el esquema
de backups a medida que crece el volumen \\
\end{longtable}

Los tres riesgos están reflejados en el supuesto de crecimiento del 10\%
anual de costos operativos usado en 5.3 y 5.5.
